\documentclass[11pt]{article}
\usepackage[margin=1in]{geometry}
\usepackage{amsmath,amssymb,amsthm,mathtools}
\usepackage{graphicx}
\usepackage[hidelinks]{hyperref}
\usepackage{microtype}
\newtheorem{theorem}{Theorem}
\newtheorem{lemma}[theorem]{Lemma}
\newcommand{\Sigmaa}{\Sigma}

\title{Operator-norm convergence of unscaled subdiagonal Pad\'e approximants\\
for bounded analytic semigroups}
\author{Candidate full solution for arXiv:1210.8408, Remark 4.7}
\date{9 August 2026}

\begin{document}
\maketitle
\begin{center}
\fbox{\parbox{0.92\textwidth}{\textbf{Status.}
Claimed full solution, likely valid, pending expert review. The result proves the
stronger operator-norm conclusion asked about in the source paper, for every bounded
analytic semigroup on a complex Banach space.}}
\end{center}

\section{Source question}
Egert and Rozendaal study the first subdiagonal Pad\'e approximants
$r_n=[n/(n+1)]$ to the exponential \cite{EgertRozendaal}. Their Theorem 4.6
proves, for a bounded analytic semigroup $T(t)=e^{-tA}$, convergence on every
fractional domain $D(A^\alpha)$, $\alpha>0$. They then ask in Remark 4.7
(arXiv PDF p.~13) whether the method \emph{without scaling and squaring}
converges strongly on the whole space and even in operator norm.
\begin{figure}[h]
  \centering
  \includegraphics[width=\textwidth]{figures/open_problem_crop.png}
  \caption{The open problem in Egert--Rozendaal, Remark 4.7, p.~13.}
\end{figure}
The answer is affirmative. In fact, the convergence below is uniform in
$t\geq0$.

\section{Notation and theorem}
Let $A$ be sectorial of angle $\varphi<\pi/2$ on a complex Banach space $X$,
and let $-A$ generate the bounded analytic semigroup $(e^{-tA})_{t\geq0}$.
For $\nu\in(\varphi,\pi/2)$ put
\[
 M_\nu:=\sup_{z\in\partial\Sigmaa_\nu}\|z(z-A)^{-1}\|,
 \qquad \delta_\nu:=1-\frac{2\nu}{\pi}>0,
\]
where $\Sigmaa_\nu=\{z\neq0:|\arg z|<\nu\}$. The Pad\'e approximant
$r_n=P_n/Q_n$ is normalized by $P_n(0)=Q_n(0)=1$, with $\deg P_n=n$ and
$\deg Q_n=n+1$, and satisfies $r_n(z)-e^z=O(z^{2n+2})$ at zero.

\begin{theorem}[Full operator-norm endpoint]\label{thm:main}
For every $\nu\in(\varphi,\pi/2)$ there is a constant $C_\nu$ such that, for
all $n\geq1$,
\[
 \sup_{t\geq0}\bigl\|r_n(-tA)-e^{-tA}\bigr\|
 \leq C_\nu M_\nu (n+1)^{-\delta_\nu}
       \bigl(1+\log(n+1)\bigr).
\]
Consequently $r_n(-tA)\to e^{-tA}$ in $\mathcal L(X)$, uniformly for
$t\geq0$. In particular, both conclusions asked for in Remark 4.7 of
\cite{EgertRozendaal} hold without any additional bounded
$H^\infty$-calculus assumption.
\end{theorem}

\section{Proof intuition}
The obstruction at the endpoint $\alpha=0$ is logarithmic: the crude
$\mathcal A$-stability estimate $|r_n(-z)-e^{-z}|\leq2$ is not integrable
against $|dz|/|z|$ on a sectorial contour. A classical theorem of
Saff--Varga--Ni supplies the missing boundary information: the unscaled
$[n/(n+1)]$ Pad\'e error is at most $1/(2(n+1))$ on the positive real axis.
The two-constants theorem transfers this decay to every strict sector, with
exponent $\delta_\nu=1-2\nu/\pi$.

Uniform sector decay alone still leaves an infinite logarithmic interval.
The second ingredient is a polynomial tail bound
$|r_n(-z)|\leq4(n+1)^2/|z|$ in the right half-plane. It follows from an exact
relation between $P_n$ and $Q_n$, followed by the Enestr\"om--Kakeya root
bound. Splitting the contour into a small piece, a middle piece, and a tail
then gives $O((n+1)^{-\delta_\nu}\log(n+1))$. The holomorphic functional
calculus turns this scalar estimate into the operator-norm estimate.

\section{Scalar estimates}
Set
\[
 N:=n+1,\qquad F_n(z):=r_n(-z)-e^{-z}
 \quad (\operatorname{Re}z\geq0).
\]

\begin{lemma}[Positive-ray and strict-sector estimates]\label{lem:sector}
For $n\geq1$,
\[
 \sup_{x\geq0}|F_n(x)|\leq\frac1{2N}.
\]
Moreover, for $0\leq\theta<\pi/2$, $r\geq0$, and
$\delta_\theta:=1-2\theta/\pi$,
\[
 |F_n(re^{\pm i\theta})|\leq 2N^{-\delta_\theta}.
\]
\end{lemma}
\begin{proof}
Write $R_{\mu,m}$ for the type $(\mu,m)$ Pad\'e approximant to $e^{-z}$ in
the notation of Saff--Varga--Ni. Their Theorem 2.1 gives
\[
 \sup_{x\geq0}|e^{-x}-R_{\mu,m}(x)|
 \leq \frac1{2^{m-\mu}\binom m\mu}
 \qquad(0\leq\mu\leq m)
\]
\cite[Theorem 2.1]{SaffVargaNi}. Taking $(\mu,m)=(N-1,N)$ gives the first
claim.

The function $F_n$ is bounded and holomorphic in the right half-plane and is
continuous on its closure. On the positive boundary ray it is bounded by
$(2N)^{-1}$, while $\mathcal A$-stability gives $|F_n(iy)|\leq2$ on the
imaginary boundary ray. The two-constants theorem in the first quadrant yields
\[
 |F_n(re^{i\theta})|
 \leq (2N)^{-\delta_\theta}2^{1-\delta_\theta}
 \leq2N^{-\delta_\theta}.
\]
The lower quadrant follows by conjugation.
\end{proof}

\begin{lemma}[Polynomial tail]\label{lem:tail}
For $n\geq1$ and $\operatorname{Re}z\geq0$, $z\neq0$,
\[
 |r_n(-z)|\leq\frac{4N^2}{|z|}.
\]
\end{lemma}
\begin{proof}
Using the explicit Pad\'e coefficients from \cite{EgertRozendaal}, write
\[
 Q_n(z)=\sum_{j=0}^{N}q_j(-z)^j,\qquad
 q_j=\frac{(2N-1-j)!N!}{(2N-1)!j!(N-j)!}>0.
\]
If $p_j$ is the coefficient of $z^j$ in $P_n$, direct cancellation gives
\[
 \frac{p_j}{q_j}=\frac{N-j}{N}\qquad(0\leq j\leq N-1).
\]
Consequently
\begin{equation}\label{eq:coefficient-identity}
 P_n(-z)=Q_n(z)-\frac zNQ_n'(z).
\end{equation}

Apply the Enestr\"om--Kakeya theorem to
$Q_n(-z)=\sum_{j=0}^{N}q_jz^j$. Every zero $\lambda$ of $Q_n$ satisfies
\[
 |\lambda|\leq\max_{0\leq j<N}\frac{q_j}{q_{j+1}}
 =\max_{0\leq j<N}\frac{(2N-1-j)(j+1)}{N-j}\leq2N^2.
\]
Let $\lambda_1,\ldots,\lambda_N$ be the roots of $Q_n$. Since $Q_n$ has
real coefficients, \eqref{eq:coefficient-identity} and logarithmic
differentiation show, for $t\in\mathbb R$, that
\begin{align*}
 |r_n(-it)|
 &=\left|1-\frac{it}{N}\frac{Q_n'(it)}{Q_n(it)}\right|\\
 &=\left|\frac1N\sum_{k=1}^{N}
             \frac{\lambda_k}{\lambda_k-it}\right|.
\end{align*}
If $|t|\leq4N^2$, $\mathcal A$-stability gives
$|t r_n(-it)|\leq4N^2$. If $|t|\geq4N^2$, then
$|\lambda_k-it|\geq|t|/2$, and the displayed formula again gives
$|t r_n(-it)|\leq4N^2$. Thus $h_n(z):=zr_n(-z)$ is bounded by $4N^2$ on
the imaginary axis. It is holomorphic on the closed right half-plane and has
a finite limit at infinity. The maximum principle on the compactified
half-plane gives $|h_n(z)|\leq4N^2$ throughout $\operatorname{Re}z\geq0$.
\end{proof}

\begin{lemma}[Contour integral]\label{lem:integral}
Fix $0<\nu<\pi/2$ and put $\delta=1-2\nu/\pi$. Then
\[
 I_{n,\nu}:=\int_0^\infty |F_n(re^{i\nu})|\,\frac{dr}{r}
 \leq C_\nu N^{-\delta}(1+\log N).
\]
The same estimate holds on the ray of angle $-\nu$.
\end{lemma}
\begin{proof}
Egert--Rozendaal's Perron estimate gives, for $\operatorname{Re}z\geq0$,
\begin{equation}\label{eq:local}
 |F_n(z)|\leq a_n|z|^{2N},\qquad
 a_n:=\frac12\left(\frac{n!}{(2n+1)!}\right)^2\leq\frac12
 \quad\cite[Lemma 3.2]{EgertRozendaal}.
\end{equation}
Let $R=N^{2+\delta}$. On $(0,1)$ use \eqref{eq:local}; on $(1,R)$ use
Lemma~\ref{lem:sector}; and on $(R,\infty)$ use Lemma~\ref{lem:tail} and
$|e^{-re^{i\nu}}|=e^{-r\cos\nu}$. This gives
\begin{align*}
 I_{n,\nu}
 &\leq \frac{a_n}{2N}+2N^{-\delta}\log R
   +\int_R^\infty
       \left(\frac{4N^2}{r}+e^{-r\cos\nu}\right)\frac{dr}{r}\\
 &\leq \frac1{4N}+2(2+\delta)N^{-\delta}\log N
   +4N^{-\delta}+\frac{e^{-R\cos\nu}}{R\cos\nu}.
\end{align*}
Since $0<\delta<1$, this is bounded by
$C_\nu N^{-\delta}(1+\log N)$. Conjugation handles the lower ray.
\end{proof}

\section{Proof of Theorem~\ref{thm:main}}
The case $t=0$ is exact. Fix $t>0$. The function
\[
 G_{n,t}(z):=r_n(-tz)-e^{-tz}=F_n(tz)
\]
belongs to the standard $H_0^\infty$ class on every sector strictly contained
in the right half-plane: it vanishes to order $2N$ at zero and is
$O(|z|^{-1})$ at infinity there. The sectorial functional calculus gives
\[
 r_n(-tA)-e^{-tA}
 =\frac1{2\pi i}\int_{\partial\Sigmaa_\nu}
      G_{n,t}(z)(z-A)^{-1}\,dz.
\]
Taking norms and using the definition of $M_\nu$ yields
\[
 \|r_n(-tA)-e^{-tA}\|
 \leq\frac{M_\nu}{2\pi}\sum_{\epsilon\in\{-1,1\}}
 \int_0^\infty |F_n(tre^{i\epsilon\nu})|\,\frac{dr}{r}.
\]
The substitution $s=tr$ removes $t$. Lemma~\ref{lem:integral} proves
\[
 \|r_n(-tA)-e^{-tA}\|
 \leq C_\nu M_\nu N^{-\delta_\nu}(1+\log N),
\]
uniformly for $t>0$, which is the assertion.

\section{Verification and limitations}
The included script \texttt{code/verify\_pade\_endpoint.py} checks the exact
coefficient identity for $1\leq n\leq20$, verifies the coefficient-ratio root
bound for $1\leq n\leq200$, and numerically samples the positive-ray,
imaginary-tail, and strict-sector integral estimates. These computations are
sanity checks only; the proof above is analytic.

The conclusion is exactly scoped to the open problem's bounded-analytic
setting. It does not settle strong convergence on all of $X$ for an arbitrary
bounded, nonanalytic $C_0$-semigroup; that broader problem is still described
as open in the 2024 work of Gomilko--Tomilov \cite{GomilkoTomilov}.

The proof deliberately uses a coarse $O(N^2/|z|)$ rational tail. Numerical
experiments suggest the sharper bound $O(N/|z|)$ and a faster contour rate,
but neither sharpening is needed for convergence and neither is claimed here.

\section{Novelty check and review recommendation}
A bounded search covered the exact Remark 4.7 wording, the core
subdiagonal-Pad\'e/operator-norm/bounded-analytic-semigroup phrases,
arXiv:1210.8408 citations, and the later papers of
Neubrander--\"Ozer--Windsperger \cite{NeubranderOzerWindsperger},
Gomilko--Tomilov \cite{GomilkoTomilov}, and
Batty--Gomilko--Tomilov \cite{BattyGomilkoTomilov}. No matching theorem for
the unscaled variable sequence $[n/(n+1)]$ on arbitrary Banach-space bounded
analytic semigroups was found. The 2024/2025 works address either the general
bounded-semigroup problem on Hilbert spaces with positive regularity, or
fixed-rational-function scaling formulas $r(tA/n)^n$.

\textbf{Human-review recommendation.} Verify first the Saff--Varga--Ni
indexing $(\mu,m)=(n,n+1)$, then \eqref{eq:coefficient-identity} and the
maximum-principle passage in Lemma~\ref{lem:tail}. If these pass, the
remaining contour-calculus argument is standard.

\begin{thebibliography}{9}
\bibitem{EgertRozendaal}
M.~Egert and J.~Rozendaal,
\emph{Convergence of subdiagonal Pad\'e approximations of $C_0$-semigroups},
J. Evol. Equ. \textbf{13} (2013), 875--895; arXiv:1210.8408.
\bibitem{SaffVargaNi}
E.~B. Saff, R.~S. Varga, and W.-C. Ni,
\emph{Geometric convergence of rational approximations to $e^{-z}$ in
infinite sectors}, Numer. Math. \textbf{26} (1976), 211--225.
\bibitem{NeubranderOzerWindsperger}
F.~Neubrander, K.~\"Ozer, and L.~Windsperger,
\emph{On subdiagonal rational Pad\'e approximations and the Brenner--Thom\'ee
approximation theorem for operator semigroups},
Discrete Contin. Dyn. Syst. S \textbf{13} (2020), 3565--3579.
\bibitem{GomilkoTomilov}
A.~Gomilko and Y.~Tomilov,
\emph{Rational approximation of operator semigroups via the
$\mathcal B$-calculus}, J. Funct. Anal. \textbf{287} (2024), 110426;
arXiv:2403.14411.
\bibitem{BattyGomilkoTomilov}
C.~Batty, A.~Gomilko, and Y.~Tomilov,
\emph{Rational approximation of holomorphic semigroups revisited},
J. Lond. Math. Soc. \textbf{111} (2025), e70066; arXiv:2403.15894.
\end{thebibliography}
\end{document}
